\section{Distinction Graph Machines}
\label{sec:main-dgm}

DGM instantiates refinement as a graph operation.
Trees impose a rooted order on distinctions; DGM stores local pairwise distinctions between runtime concepts and keeps repairable memories at the nodes.
This is useful when conflicts are local and non-hierarchical: a new concept may need to be separated from only the concepts it has been confused with, rather than from every existing leaf in a tree.
The runtime state is
\[
S_t=
\left(
\{a_j,c_j,M_j,n_j\}_{j=1}^{N_t},
E_t
\right).
\]
Here $a_j\in\R^d$ is a frozen anchor that defines the hard information structure, $c_j\in\R^d$ is a movable centroid used only for retrieval, $M_j$ is a repairable local memory, and $E_t$ is the set of accepted distinction edges.
Anchors and edges define refinement; memories and centroids are repairable payloads and statistics.
A new concept is initialized with anchor and centroid at the current hidden state; later repair may move the centroid or update the memory, but does not rewrite previously accepted edge predicates.

For an edge $e=(i,j)$,
\[
u_e=\frac{a_i-a_j}{\|a_i-a_j\|_2},
\qquad
b_e=\frac12u_e^\top(a_i+a_j),
\qquad
g_e(h)=\mathbf 1\{u_e^\top h>b_e\}.
\]
This edge separates the anchors that caused the conflict:
\[
u_e^\top a_i-b_e=\frac12\|a_i-a_j\|_2,
\qquad
u_e^\top a_j-b_e=-\frac12\|a_i-a_j\|_2.
\]

The edge signature
\[
\sigma^G_t(h)=\bigl(g_e(h)\bigr)_{e\in E_t}
\]
induces the edge partition
\[
\calP^G_t
=
\{h:\sigma^G_t(h)=s\}_{s\in\{0,1\}^{E_t}}
\setminus\{\varnothing\}.
\]
Therefore
\[
\mathrm{Repair}\Rightarrow \calP^G_{t+1}=\calP^G_t,
\qquad
\mathrm{Refine\ by\ }g\Rightarrow \calP^G_{t+1}=\calP^G_t\vee g.
\]
These are the DGM invariants.
A repair may improve $M_j$, change counts, or move retrieval centroids, but it does not add an edge generator.
A refinement creates a new anchor and at least one accepted edge, so the edge signature gains a new bit and the edge skeleton is joined with that bit.
The exact hard route also uses anchor-order events for nearest eligible anchor selection; the full formalism is in Appendix~\ref{sec:algorithms}.
The edge skeleton is the part that corresponds exactly to the minimal join operation in Section~\ref{sec:main-value}.
Anchor-order events are routing infrastructure, not additional consequence claims.

The practical Soft-DGM read is a sparse differentiable approximation.
Given
\[
N(h)=
N_{\rm centroid}(h)
\cup
N_{\rm anchor}(h)
\cup
\operatorname{GraphNeighbors}(N_{\rm centroid}(h)),
\]
define
\[
\chi_j(h)
=
\exp\left[
\frac1{\sqrt{\max\{1,|\operatorname{Inc}(j)|\}}}
\sum_{e\in\operatorname{Inc}(j)}
\beta_e
\log\operatorname{sigmoid}
\left(
\frac{s_{e,j}(u_e^\top h-b_e)}{\tau_e}
\right)
\right].
\]
The routing weight and prediction are
\[
\alpha_j(h)
=
\frac{
\exp(-\|h-c_j\|_2^2/T)\chi_j(h)
}{
\sum_{r\in N(h)}
\exp(-\|h-c_r\|_2^2/T)\chi_r(h)
},
\qquad
\hat y(h)=\sum_{j\in N(h)}\alpha_j(h)M_j(h).
\]
Exact partition guarantees apply to the anchor-based hard route.
The soft read is a sparse differentiable approximation and requires candidate recall; diagnostics and the formal approximation statement are in Appendix~\ref{sec:algorithms}.
Centroid movement can change soft routing weights, so centroid repair should be interpreted as repair of the hard information structure plus an approximate neural read, not as a claim that every smooth prediction remains constant on edge-partition atoms.

\begin{algorithm}[t]
\caption{\textsc{Guarded-DGM-Update}}
\label{alg:main-guarded-dgm}
\begin{algorithmic}[1]
\Statex Input: state $S_t$, example $(x_t,y_t)$, proposal buffer, scoring buffer.
\State Compute $h_t=\phi_\theta(x_t)$ and predict with edge-gated local memories.
\State Form $\calC_t^{\rm rep}$ by updating local memory or centroid while keeping anchors and edges fixed.
\State Form $\calC_t^{\rm ref}$ by creating a new anchor at $h_t$ and edges to confused neighbors.
\State Score all candidates on $\calB_t^{\rm score}$.
\State Let $G_{\rm rep}$ be the best empirical repair gain.
\State Let $G_{\rm ref}$ be the best penalized empirical refinement gain.
\If{$G_{\rm ref}>G_{\rm rep}+\tau+4\epsilon_t$}
    \State commit the refinement.
\Else
    \State commit the best non-worsening repair candidate.
\EndIf
\end{algorithmic}
\end{algorithm}

DGM is therefore a data-structure implementation of guarded distinction acceptance under an explicit runtime information structure.
It is not a claim that graph memory is globally optimal, nor that the soft implementation is exactly measurable with respect to the hard edge partition.
The claim boundary is important: Hard-DGM gives the clean partition invariant, Soft-DGM gives a practical sparse read, and the guard decides whether the empirical evidence justifies moving from repair to refinement.
Thus DGM should be evaluated by whether it makes the right structural decision under a stated runtime contract, not by whether its graph edges are human-interpretable concepts by default.
